\section{Native Defense Audit Methods}
\label{sec:native-audit-methods}

We evaluate whether published gradient-leakage defenses retain their claimed
protection when assessed in their native reconstruction settings and under an
attacker whose knowledge matches the server threat model.  The audit is
pre-registered in the Priority 29 and Priority 30 amendments.  It treats a
defense as an empirical mechanism, rather than inferring privacy from a
transformed gradient alone.  All emitted cells use fresh, source-disjoint
targets, single-threaded execution, and per-target result records.

\subsection{Native instruments and positive controls}
For images, the native instrument is CIFAR-10 test data, an untrained official
\texttt{LeNetZhu} model constructed with seed 42, FedSGD with one known-label
example, and the official Geiping \texttt{GradientReconstructor}.  The attack
uses cosine cost, signed Adam with learning rate 0.1, 4,800 iterations, one
restart, total variation weight 0.01, boxed image projection, learning-rate
decay, and attacker-observable loss for restart scoring.  We score raw-space
MSE, PSNR, and SSIM against the target, and retain gray 0.5 and CIFAR training
mean image baselines.

For tabular data, the native instrument is the official TabLeak Adult FedSGD
configuration: known labels; batch size 8; an initialized fully connected
network with dimensions $105\rightarrow100\rightarrow100\rightarrow2$;
cross-entropy; cosine discrepancy; signed Adam with learning rate 0.06; 1,500
iterations; 30 restart/ensemble members; uniform initialization; categorical
softmax; bounded continuous sigmoid parameterization; median-plus-softmax
pooling; and no learning-rate schedule.  The primary tabular outcome is the
official mixed-feature reconstruction accuracy, with mean/mode and empirical
marginal baselines.

Before defense comparisons, each native instrument must pass a positive-control
gate at $n=8$ and then $n=24$ fresh targets.  The image attack must outperform
the Prior, decoy, gray, and CIFAR-mean baselines; the Adult attack must
outperform both data-free tabular baselines.  Each directional gate uses an
exact one-sided sign test at $\alpha=0.05$.  A failed native gate stops the
domain and produces \texttt{NOT\_ASSESSABLE}, rather than evidence for a
defense.

\subsection{Defense adapters and checks}
For each defense we first require four checks: (i) defense-not-identity at the
frozen published hyperparameters; (ii) lossless sanity for an adaptive path
when such a setting exists; (iii) an explicit server-knowledge receipt for
every key, hash, mask, stochastic sample, or metadata needed to decode the
payload; and (iv) decoy correctness plus the native positive-control gate.
PRECODE is the official variational bottleneck inserted immediately before the
final classifier and evaluated through its stochastic forward pass.  Soteria
uses client-side representation sensitivity and masks final-classifier gradient
columns at 80\% for images and 40\% for Adult.  Gradient pruning applies
per-tensor lowest-magnitude pruning at 70\%.  ATS is image-only and applies the
official searched policy 3-1-7 before gradient computation, while scoring both
against the raw image and, for the paper comparison, the transformed image.
Count-Sketch uses five rows, 500,000 columns, seed 21, and registered
trainable-gradient tensor order; known- and unknown-hash server conditions are
distinct.  Unknown hashes make an adaptive decoder unavailable rather than
permitting an invented sketch operator.

\subsection{Evaluation matrix and adaptive attackers}
Each qualified defense-domain cell has $n=39$ confirmatory targets.  E1 uses
the non-adaptive official attacker and the original paper's metric, asking
whether the published qualitative protection claim reproduces.  E2 uses an
adaptive attacker matched to the declared server knowledge and tests whether
reconstruction quality is reduced versus the undefended attack on the same
targets.  The PRECODE attacker differentiates through the stochastic
bottleneck; the Soteria attacker receives the representation mask; the pruning
attacker uses the retained-coordinate observation; ATS uses the known input
transformation; and known-hash Count-Sketch uses the actual linear sketch map.
For DNA v2, the key-holder path uses least squares through the sketch map, not
a naive lift.  For every unknown-key or unknown-hash condition, unavailable
knowledge is recorded as unavailable rather than assumed.

E3 provides three comparators: (i) paper-documented unclipped Gaussian noise
at the paper's noise scale, (ii) properly clipped DP at the defense's matched
$L_2$ distortion, and (iii), for tabular cells only, utility-matched DP by
Adult FedSGD test accuracy using the pre-specified bracketed-grid rule.  E3(iii)
is run only if an E2 tabular defense survives and the frozen wall-clock rule
permits it.  The frozen reduction order may also omit ATS and the secondary
PaySim reconstruction analysis; these omissions are reported as such, not
treated as negative results.

For every paired comparison we use exact one-sided sign tests, excluding ties,
at $n=39$.  Holm correction is applied over the whole Priority 29 family of E2
and E3 p-values.  \todo{Insert the emitted comparison family, raw p-values,
and Holm-adjusted p-values.} A defense \texttt{SURVIVES} only if E2 shows a
significant reconstruction-quality reduction versus no defense and is not
significantly worse than properly clipped matched-distortion DP (and, for
Adult, matched-utility DP).  It \texttt{DOES\_NOT\_SURVIVE} otherwise, with
the failing check named; it is \texttt{NOT\_ASSESSABLE} if a native positive
control fails or a required model-level adapter cannot be supplied without
altering read-only reference code.  \todo{Insert per-defense RQ4 verdicts
only after the paired comparator records and corrected tests are available.}
